\documentclass{beamer}
% =====================================================================
% finance-beamer.tex — shared Beamer style for the LLM-in-Finance course
% =====================================================================
% Reusable slide style. Each deck \input's this immediately after
% \documentclass{beamer}. It provides:
%   * the Madrid theme + book colour palette (matches the printed book)
%   * two book-style framing blocks:
%       \begin{bigpicture} ... \end{bigpicture}   ("the bigger picture")
%       \begin{underhood}   ... \end{underhood}    ("under the hood")
%     These mirror the book's `context` and `deepdive` tcolorboxes so the
%     same intuition-vs-mechanism scaffold the reader meets in the book
%     reappears on the slides.
%   * a Python listings style for code frames
%   * appendix conventions: \appendixstart drops a divider and switches
%     to non-numbered backup frames so "extra / less central" material
%     lives after the main talk without inflating the slide count.
%
% Usage:
%   \documentclass{beamer}
%   % =====================================================================
% finance-beamer.tex — shared Beamer style for the LLM-in-Finance course
% =====================================================================
% Reusable slide style. Each deck \input's this immediately after
% \documentclass{beamer}. It provides:
%   * the Madrid theme + book colour palette (matches the printed book)
%   * two book-style framing blocks:
%       \begin{bigpicture} ... \end{bigpicture}   ("the bigger picture")
%       \begin{underhood}   ... \end{underhood}    ("under the hood")
%     These mirror the book's `context` and `deepdive` tcolorboxes so the
%     same intuition-vs-mechanism scaffold the reader meets in the book
%     reappears on the slides.
%   * a Python listings style for code frames
%   * appendix conventions: \appendixstart drops a divider and switches
%     to non-numbered backup frames so "extra / less central" material
%     lives after the main talk without inflating the slide count.
%
% Usage:
%   \documentclass{beamer}
%   % =====================================================================
% finance-beamer.tex — shared Beamer style for the LLM-in-Finance course
% =====================================================================
% Reusable slide style. Each deck \input's this immediately after
% \documentclass{beamer}. It provides:
%   * the Madrid theme + book colour palette (matches the printed book)
%   * two book-style framing blocks:
%       \begin{bigpicture} ... \end{bigpicture}   ("the bigger picture")
%       \begin{underhood}   ... \end{underhood}    ("under the hood")
%     These mirror the book's `context` and `deepdive` tcolorboxes so the
%     same intuition-vs-mechanism scaffold the reader meets in the book
%     reappears on the slides.
%   * a Python listings style for code frames
%   * appendix conventions: \appendixstart drops a divider and switches
%     to non-numbered backup frames so "extra / less central" material
%     lives after the main talk without inflating the slide count.
%
% Usage:
%   \documentclass{beamer}
%   \input{../_style/finance-beamer.tex}
%   \title{...}\subtitle{...}\author{...}\institute{...}\date{\today}
%   \begin{document} ... \end{document}
% =====================================================================

\usetheme{Madrid}
\usecolortheme{whale}
\setbeamertemplate{navigation symbols}{}
\setbeamertemplate{caption}[numbered]
\setbeamercovered{transparent}

% --- Density controls: keep dense, equation-heavy frames on one slide ---
% Slightly smaller body text + tighter lists/blocks. Frame titles and the
% Madrid chrome keep their normal size, so the deck still reads as Madrid,
% just with more vertical room for content.
\setbeamerfont{normal text}{size=\small}
\setbeamertemplate{itemize/enumerate body begin}{\small}
\setbeamertemplate{itemize/enumerate subbody begin}{\footnotesize}
% Tighter block spacing.
\setbeamertemplate{block begin}{%
  \par\vskip2pt%
  \begin{beamercolorbox}[colsep*=.75ex,leftskip=1ex,rightskip=1ex]{block title}%
    \usebeamerfont*{block title}\insertblocktitle%
  \end{beamercolorbox}%
  {\parskip0pt\vskip-1pt}%
  \begin{beamercolorbox}[colsep*=.75ex,vmode,leftskip=1ex,rightskip=1ex]{block body}%
  \vskip1pt}
\setbeamertemplate{block end}{\end{beamercolorbox}\vskip2pt}

\usepackage{amsmath, amssymb}
\usepackage{booktabs}
\usepackage{xcolor}
\usepackage{listings}
\usepackage[skins, breakable]{tcolorbox}

% --- Book colour palette (kept in sync with book/preamble.tex) ---
\definecolor{linkblue}{RGB}{14, 75, 140}
\definecolor{deepdivebg}{RGB}{235, 244, 255}
\definecolor{narrativebg}{RGB}{255, 252, 240}
\definecolor{narrativerule}{RGB}{180, 130, 40}
\definecolor{steelblue}{RGB}{70, 130, 180}
\definecolor{forestgreen}{RGB}{34, 139, 34}
\definecolor{crimson}{RGB}{220, 20, 60}
\definecolor{darkorange}{RGB}{255, 140, 0}
\definecolor{codebg}{RGB}{248, 248, 248}
\definecolor{coderule}{RGB}{210, 210, 210}
\definecolor{keywordblue}{RGB}{0, 100, 180}
\definecolor{commentgray}{RGB}{120, 120, 120}
\definecolor{stringgreen}{RGB}{30, 130, 30}

% Tie the Madrid structural colour to the book's link blue.
\setbeamercolor{structure}{fg=linkblue}
\setbeamercolor{title}{fg=white}
\setbeamercolor{frametitle}{fg=white}

% --- "the bigger picture" : narrative / intuition framing -----------
\newtcolorbox{bigpicture}{%
  enhanced, breakable, boxsep=2pt,
  colback  = narrativebg,
  colframe = narrativerule,
  boxrule  = 0pt, toprule = 1.4pt, bottomrule = 0.4pt,
  leftrule = 0pt, rightrule = 0pt, arc = 0pt,
  left = 6pt, right = 6pt, top = 4pt, bottom = 5pt,
  fonttitle = \footnotesize\scshape\color{narrativerule},
  title = {the bigger picture}, attach title to upper = {\par\smallskip},
}

% --- "under the hood" : the mechanism / technical detail ------------
\newtcolorbox{underhood}{%
  enhanced, breakable, boxsep=2pt,
  colback  = deepdivebg,
  colframe = linkblue,
  boxrule  = 0pt, toprule = 1.4pt, bottomrule = 0.4pt,
  leftrule = 0pt, rightrule = 0pt, arc = 0pt,
  left = 6pt, right = 6pt, top = 4pt, bottom = 5pt,
  fonttitle = \footnotesize\scshape\color{linkblue},
  title = {under the hood}, attach title to upper = {\par\smallskip},
}

% --- Python code style ---------------------------------------------
\lstset{
  basicstyle    = \ttfamily\footnotesize,
  keywordstyle  = \color{keywordblue}\bfseries,
  commentstyle  = \color{commentgray}\itshape,
  stringstyle   = \color{stringgreen},
  showstringspaces = false,
  language      = Python,
  breaklines    = true,
  backgroundcolor = \color{codebg},
  frame         = single,
  rulecolor     = \color{coderule},
  numbers       = none,
  aboveskip     = 4pt, belowskip = 4pt,
}

% --- Appendix conventions ------------------------------------------
% \appendixstart{Title} : begins the "extra material" appendix. Frames
% after it are not counted toward the main deck's frame total, keeping
% the core talk lean while preserving the deeper / optional content.
\newcommand{\appendixstart}[1]{%
  \appendix
  \setbeamertemplate{frametitle continuation}{}%
  \begin{frame}[noframenumbering]
    \begin{center}
      {\Large\scshape\color{linkblue} Appendix}\\[4pt]
      {\large #1}\\[10pt]
      {\footnotesize Extra / optional material — beyond the core lecture.}
    \end{center}
  \end{frame}
}
% Use \begin{frame}[noframenumbering]{...} for each appendix frame.

%   \title{...}\subtitle{...}\author{...}\institute{...}\date{\today}
%   \begin{document} ... \end{document}
% =====================================================================

\usetheme{Madrid}
\usecolortheme{whale}
\setbeamertemplate{navigation symbols}{}
\setbeamertemplate{caption}[numbered]
\setbeamercovered{transparent}

% --- Density controls: keep dense, equation-heavy frames on one slide ---
% Slightly smaller body text + tighter lists/blocks. Frame titles and the
% Madrid chrome keep their normal size, so the deck still reads as Madrid,
% just with more vertical room for content.
\setbeamerfont{normal text}{size=\small}
\setbeamertemplate{itemize/enumerate body begin}{\small}
\setbeamertemplate{itemize/enumerate subbody begin}{\footnotesize}
% Tighter block spacing.
\setbeamertemplate{block begin}{%
  \par\vskip2pt%
  \begin{beamercolorbox}[colsep*=.75ex,leftskip=1ex,rightskip=1ex]{block title}%
    \usebeamerfont*{block title}\insertblocktitle%
  \end{beamercolorbox}%
  {\parskip0pt\vskip-1pt}%
  \begin{beamercolorbox}[colsep*=.75ex,vmode,leftskip=1ex,rightskip=1ex]{block body}%
  \vskip1pt}
\setbeamertemplate{block end}{\end{beamercolorbox}\vskip2pt}

\usepackage{amsmath, amssymb}
\usepackage{booktabs}
\usepackage{xcolor}
\usepackage{listings}
\usepackage[skins, breakable]{tcolorbox}

% --- Book colour palette (kept in sync with book/preamble.tex) ---
\definecolor{linkblue}{RGB}{14, 75, 140}
\definecolor{deepdivebg}{RGB}{235, 244, 255}
\definecolor{narrativebg}{RGB}{255, 252, 240}
\definecolor{narrativerule}{RGB}{180, 130, 40}
\definecolor{steelblue}{RGB}{70, 130, 180}
\definecolor{forestgreen}{RGB}{34, 139, 34}
\definecolor{crimson}{RGB}{220, 20, 60}
\definecolor{darkorange}{RGB}{255, 140, 0}
\definecolor{codebg}{RGB}{248, 248, 248}
\definecolor{coderule}{RGB}{210, 210, 210}
\definecolor{keywordblue}{RGB}{0, 100, 180}
\definecolor{commentgray}{RGB}{120, 120, 120}
\definecolor{stringgreen}{RGB}{30, 130, 30}

% Tie the Madrid structural colour to the book's link blue.
\setbeamercolor{structure}{fg=linkblue}
\setbeamercolor{title}{fg=white}
\setbeamercolor{frametitle}{fg=white}

% --- "the bigger picture" : narrative / intuition framing -----------
\newtcolorbox{bigpicture}{%
  enhanced, breakable, boxsep=2pt,
  colback  = narrativebg,
  colframe = narrativerule,
  boxrule  = 0pt, toprule = 1.4pt, bottomrule = 0.4pt,
  leftrule = 0pt, rightrule = 0pt, arc = 0pt,
  left = 6pt, right = 6pt, top = 4pt, bottom = 5pt,
  fonttitle = \footnotesize\scshape\color{narrativerule},
  title = {the bigger picture}, attach title to upper = {\par\smallskip},
}

% --- "under the hood" : the mechanism / technical detail ------------
\newtcolorbox{underhood}{%
  enhanced, breakable, boxsep=2pt,
  colback  = deepdivebg,
  colframe = linkblue,
  boxrule  = 0pt, toprule = 1.4pt, bottomrule = 0.4pt,
  leftrule = 0pt, rightrule = 0pt, arc = 0pt,
  left = 6pt, right = 6pt, top = 4pt, bottom = 5pt,
  fonttitle = \footnotesize\scshape\color{linkblue},
  title = {under the hood}, attach title to upper = {\par\smallskip},
}

% --- Python code style ---------------------------------------------
\lstset{
  basicstyle    = \ttfamily\footnotesize,
  keywordstyle  = \color{keywordblue}\bfseries,
  commentstyle  = \color{commentgray}\itshape,
  stringstyle   = \color{stringgreen},
  showstringspaces = false,
  language      = Python,
  breaklines    = true,
  backgroundcolor = \color{codebg},
  frame         = single,
  rulecolor     = \color{coderule},
  numbers       = none,
  aboveskip     = 4pt, belowskip = 4pt,
}

% --- Appendix conventions ------------------------------------------
% \appendixstart{Title} : begins the "extra material" appendix. Frames
% after it are not counted toward the main deck's frame total, keeping
% the core talk lean while preserving the deeper / optional content.
\newcommand{\appendixstart}[1]{%
  \appendix
  \setbeamertemplate{frametitle continuation}{}%
  \begin{frame}[noframenumbering]
    \begin{center}
      {\Large\scshape\color{linkblue} Appendix}\\[4pt]
      {\large #1}\\[10pt]
      {\footnotesize Extra / optional material — beyond the core lecture.}
    \end{center}
  \end{frame}
}
% Use \begin{frame}[noframenumbering]{...} for each appendix frame.

%   \title{...}\subtitle{...}\author{...}\institute{...}\date{\today}
%   \begin{document} ... \end{document}
% =====================================================================

\usetheme{Madrid}
\usecolortheme{whale}
\setbeamertemplate{navigation symbols}{}
\setbeamertemplate{caption}[numbered]
\setbeamercovered{transparent}

% --- Density controls: keep dense, equation-heavy frames on one slide ---
% Slightly smaller body text + tighter lists/blocks. Frame titles and the
% Madrid chrome keep their normal size, so the deck still reads as Madrid,
% just with more vertical room for content.
\setbeamerfont{normal text}{size=\small}
\setbeamertemplate{itemize/enumerate body begin}{\small}
\setbeamertemplate{itemize/enumerate subbody begin}{\footnotesize}
% Tighter block spacing.
\setbeamertemplate{block begin}{%
  \par\vskip2pt%
  \begin{beamercolorbox}[colsep*=.75ex,leftskip=1ex,rightskip=1ex]{block title}%
    \usebeamerfont*{block title}\insertblocktitle%
  \end{beamercolorbox}%
  {\parskip0pt\vskip-1pt}%
  \begin{beamercolorbox}[colsep*=.75ex,vmode,leftskip=1ex,rightskip=1ex]{block body}%
  \vskip1pt}
\setbeamertemplate{block end}{\end{beamercolorbox}\vskip2pt}

\usepackage{amsmath, amssymb}
\usepackage{booktabs}
\usepackage{xcolor}
\usepackage{listings}
\usepackage[skins, breakable]{tcolorbox}

% --- Book colour palette (kept in sync with book/preamble.tex) ---
\definecolor{linkblue}{RGB}{14, 75, 140}
\definecolor{deepdivebg}{RGB}{235, 244, 255}
\definecolor{narrativebg}{RGB}{255, 252, 240}
\definecolor{narrativerule}{RGB}{180, 130, 40}
\definecolor{steelblue}{RGB}{70, 130, 180}
\definecolor{forestgreen}{RGB}{34, 139, 34}
\definecolor{crimson}{RGB}{220, 20, 60}
\definecolor{darkorange}{RGB}{255, 140, 0}
\definecolor{codebg}{RGB}{248, 248, 248}
\definecolor{coderule}{RGB}{210, 210, 210}
\definecolor{keywordblue}{RGB}{0, 100, 180}
\definecolor{commentgray}{RGB}{120, 120, 120}
\definecolor{stringgreen}{RGB}{30, 130, 30}

% Tie the Madrid structural colour to the book's link blue.
\setbeamercolor{structure}{fg=linkblue}
\setbeamercolor{title}{fg=white}
\setbeamercolor{frametitle}{fg=white}

% --- "the bigger picture" : narrative / intuition framing -----------
\newtcolorbox{bigpicture}{%
  enhanced, breakable, boxsep=2pt,
  colback  = narrativebg,
  colframe = narrativerule,
  boxrule  = 0pt, toprule = 1.4pt, bottomrule = 0.4pt,
  leftrule = 0pt, rightrule = 0pt, arc = 0pt,
  left = 6pt, right = 6pt, top = 4pt, bottom = 5pt,
  fonttitle = \footnotesize\scshape\color{narrativerule},
  title = {the bigger picture}, attach title to upper = {\par\smallskip},
}

% --- "under the hood" : the mechanism / technical detail ------------
\newtcolorbox{underhood}{%
  enhanced, breakable, boxsep=2pt,
  colback  = deepdivebg,
  colframe = linkblue,
  boxrule  = 0pt, toprule = 1.4pt, bottomrule = 0.4pt,
  leftrule = 0pt, rightrule = 0pt, arc = 0pt,
  left = 6pt, right = 6pt, top = 4pt, bottom = 5pt,
  fonttitle = \footnotesize\scshape\color{linkblue},
  title = {under the hood}, attach title to upper = {\par\smallskip},
}

% --- Python code style ---------------------------------------------
\lstset{
  basicstyle    = \ttfamily\footnotesize,
  keywordstyle  = \color{keywordblue}\bfseries,
  commentstyle  = \color{commentgray}\itshape,
  stringstyle   = \color{stringgreen},
  showstringspaces = false,
  language      = Python,
  breaklines    = true,
  backgroundcolor = \color{codebg},
  frame         = single,
  rulecolor     = \color{coderule},
  numbers       = none,
  aboveskip     = 4pt, belowskip = 4pt,
}

% --- Appendix conventions ------------------------------------------
% \appendixstart{Title} : begins the "extra material" appendix. Frames
% after it are not counted toward the main deck's frame total, keeping
% the core talk lean while preserving the deeper / optional content.
\newcommand{\appendixstart}[1]{%
  \appendix
  \setbeamertemplate{frametitle continuation}{}%
  \begin{frame}[noframenumbering]
    \begin{center}
      {\Large\scshape\color{linkblue} Appendix}\\[4pt]
      {\large #1}\\[10pt]
      {\footnotesize Extra / optional material — beyond the core lecture.}
    \end{center}
  \end{frame}
}
% Use \begin{frame}[noframenumbering]{...} for each appendix frame.


\title{Practical Session 1: From Text to Vectors}
\subtitle{Hands-On: TF-IDF, Cosine Similarity, and Word Embeddings}
\author{Juan F. Imbet}
\institute{Paris Dauphine -- PSL University}
\date{\today}

\begin{document}

\begin{frame}
  \titlepage
\end{frame}

% ---------------------------------------------------------------
\begin{frame}{Session Overview}
  \textbf{Duration:} 1 hour \quad \textbf{Format:} guided problems + Python case study
  \medskip
  \begin{enumerate}
    \item \textbf{Recap} (5 min) --- key formulas from the lecture
    \item \textbf{Problem 1} (15 min) --- TF-IDF by hand on earnings call snippets
    \item \textbf{Problem 2} (10 min) --- interpreting word embedding analogies
    \item \textbf{Case Study} (20 min) --- Python: TF-IDF and cosine similarity on real financial text
    \item \textbf{Discussion} (5 min) --- when BoW vs.\ embeddings?
    \item \textbf{Solutions} (5 min) --- walkthrough
  \end{enumerate}
  \medskip
  Notebook: \texttt{code/notebooks/01-intro/practical.ipynb}
\end{frame}

% ---------------------------------------------------------------
\section{Recap}

\begin{frame}{Recap: TF-IDF Formulas}
  \[
    \mathrm{tf}(v, d) = \frac{\mathrm{count}(v, d)}{|d|}
    \qquad
    \mathrm{idf}(v, \mathcal{C}) = \log\!\left(\frac{|\mathcal{C}|}{\mathrm{df}(v, \mathcal{C})}\right)
  \]
  \[
    \mathrm{tfidf}(v, d, \mathcal{C}) = \mathrm{tf}(v, d) \times \mathrm{idf}(v, \mathcal{C})
  \]
  \medskip
  Document similarity (cosine, length-invariant):
  \[
    \mathrm{sim}(d_i, d_j)
    = \frac{\tilde{\mathbf{x}}(d_i)^\top \tilde{\mathbf{x}}(d_j)}
           {\|\tilde{\mathbf{x}}(d_i)\|_2\,\|\tilde{\mathbf{x}}(d_j)\|_2}
  \]
  \medskip
  \textbf{Key intuitions:}
  \begin{itemize}
    \item Word in \emph{all} documents $\Rightarrow$ IDF $= 0$ (suppressed)
    \item Word in \emph{one} document $\Rightarrow$ IDF $= \ln N$ (maximally discriminative)
    \item TF-IDF picks out words that are \emph{frequent in a document but rare in the corpus}
  \end{itemize}
\end{frame}

\begin{frame}{Recap: Word2Vec Skip-Gram with Negative Sampling}
  Skip-Gram objective (predict context from centre):
  \[
    J_{\mathrm{SG}} = \frac{1}{T}\sum_{t}\sum_{k \neq 0}^{\pm c}
    \log P(w_{t+k} \mid w_t)
  \]
  Negative sampling approximation:
  \[
    \mathcal{L}_{\mathrm{NEG}}(w_t, w_{t+k})
    = \log \sigma(\mathbf{v}_{w_{t+k}}^\top \hat{\mathbf{v}}_{w_t})
    + \sum_{j=1}^{K} \mathbb{E}_{w_j \sim P_n}
      \bigl[\log \sigma(-\mathbf{v}_{w_j}^\top \hat{\mathbf{v}}_{w_t})\bigr]
  \]
  \medskip
  \textbf{Key properties:}
  \begin{itemize}
    \item Only $K{+}1$ embedding rows updated per step (vs.\ $V$ for full softmax)
    \item Positive context vectors pushed \emph{toward} centre; negative samples pushed \emph{away}
    \item Linear analogies emerge: $\hat{\mathbf{v}}_{\text{king}} - \hat{\mathbf{v}}_{\text{man}}
          + \hat{\mathbf{v}}_{\text{woman}} \approx \hat{\mathbf{v}}_{\text{queen}}$
  \end{itemize}
\end{frame}

% ---------------------------------------------------------------
\section{Problem 1: TF-IDF by Hand}

\begin{frame}{Problem 1: Setup}
  Consider the following four earnings call snippets.
  After lower-casing and removing stop words \{the, a, of, in, by, is, our, we, to, and\}:
  \medskip
  \begin{center}
  \begin{tabular}{@{}cl@{}}
    \toprule
    $d$ & Pre-processed tokens \\
    \midrule
    $d_1$ & \texttt{revenue increased strong margin growth} \\
    $d_2$ & \texttt{margin pressure volume declined uncertainty} \\
    $d_3$ & \texttt{strong volume revenue outlook positive} \\
    $d_4$ & \texttt{guidance withdrawn uncertainty pressure growth} \\
    \bottomrule
  \end{tabular}
  \end{center}
  \medskip
  Vocabulary ($V = 10$): \texttt{declined, growth, guidance, increased, margin, outlook,
  positive, pressure, revenue, strong, uncertainty, volume, withdrawn}\\
  (remove duplicates; $V = 13$)
\end{frame}

\begin{frame}{Problem 1: Tasks}
  Working individually or in pairs:
  \medskip
  \begin{enumerate}
    \item \textbf{Compute} $\mathrm{tf}(v, d)$ for each token $v$ and document $d$.
          \textit{Hint:} each document has 5 tokens, so $\mathrm{tf} \in \{0, 0.2\}$.
    \medskip
    \item \textbf{Compute} $\mathrm{df}(v, \mathcal{C})$ and
          $\mathrm{idf}(v, \mathcal{C}) = \ln(4 / \mathrm{df}(v))$ for each vocabulary word.
    \medskip
    \item \textbf{Compute} $\mathrm{tfidf}(v, d_1, \mathcal{C})$ and
          $\mathrm{tfidf}(v, d_2, \mathcal{C})$ for all $v$.
    \medskip
    \item \textbf{Compute} $\mathrm{sim}(d_1, d_2)$ using $\ell_2$-normalised TF-IDF vectors.
          \textit{Hint:} only non-zero entries contribute to the dot product.
    \medskip
    \item \textbf{Interpret:} which two documents are most similar? Does this match
          your intuition about their content?
  \end{enumerate}
  \medskip
  \small Time: 15 minutes.
\end{frame}

% ---------------------------------------------------------------
\section{Problem 2: Word Embedding Analogies}

\begin{frame}{Problem 2: Setup and Tasks}
  Suppose you have a Word2Vec model trained on a large corpus of earnings call transcripts.
  The model has learned the following (approximate) relationships in embedding space:

  \begin{center}
  \begin{tabular}{@{}ll@{}}
    \toprule
    Analogy & Relationship \\
    \midrule
    \texttt{bond} $-$ \texttt{coupon} $+$ \texttt{dividend} $\approx$ ? & ? \\
    \texttt{long} $-$ \texttt{profit} $+$ \texttt{loss} $\approx$ ? & ? \\
    \texttt{inflation} $-$ \texttt{high} $+$ \texttt{low} $\approx$ ? & ? \\
    \bottomrule
  \end{tabular}
  \end{center}
  \medskip
  \textbf{Tasks:}
  \begin{enumerate}
    \item For each row, predict the missing word and explain the financial reasoning.
    \item Explain \emph{why} these linear relationships emerge from the Skip-Gram training objective.
    \item Give an example of an analogy that would \emph{fail} or produce a nonsensical result,
          and explain why.
  \end{enumerate}
  \medskip
  \small Time: 10 minutes. Discuss with a neighbour.
\end{frame}

% ---------------------------------------------------------------
\section{Case Study: Python on Earnings Calls}

\begin{frame}{Case Study: TF-IDF on Earnings Call Snippets}
  \textbf{Goal:} Use \texttt{scikit-learn} to compute TF-IDF representations
  of five real earnings call snippets and identify the most similar pair.
  \medskip
  \textbf{Setup:}
  \begin{itemize}
    \item Five snippets from Q3 2023 earnings calls (provided in the notebook)
    \item Apply \texttt{TfidfVectorizer} with \texttt{stop\_words='english'}, \texttt{max\_features=500}
    \item Compute pairwise cosine similarity matrix
    \item Identify and interpret the most similar and least similar pairs
  \end{itemize}
  \medskip
  Open \texttt{code/notebooks/01-intro/practical.ipynb} and run cells 1--4.
\end{frame}

\begin{frame}[fragile]{Case Study: Python Code Skeleton}
\begin{lstlisting}
from sklearn.feature_extraction.text import TfidfVectorizer
from sklearn.metrics.pairwise import cosine_similarity
import numpy as np

# snippets: list of 5 earnings call excerpt strings
vectorizer = TfidfVectorizer(stop_words='english',
                              max_features=500)
X = vectorizer.fit_transform(snippets)  # (5, 500) sparse matrix

# Pairwise cosine similarity
sim_matrix = cosine_similarity(X)
np.fill_diagonal(sim_matrix, 0)  # zero out self-similarity

# Most similar pair
i, j = np.unravel_index(sim_matrix.argmax(), sim_matrix.shape)
print(f"Most similar: doc {i} and doc {j}, sim = {sim_matrix[i,j]:.3f}")

# Top terms for doc i
feature_names = vectorizer.get_feature_names_out()
top_terms = feature_names[X[i].toarray().argsort()[0, -5:]]
print(f"Top terms for doc {i}: {top_terms}")
\end{lstlisting}
\end{frame}

\begin{frame}{Case Study: Questions to Answer}
  After running the notebook:
  \medskip
  \begin{enumerate}
    \item Which pair of documents has the highest cosine similarity?
          What do their top TF-IDF terms have in common?
    \medskip
    \item The similarity between two documents from \emph{different} companies
          turns out to be higher than the similarity between two documents from the
          \emph{same} company but in different sectors. How can you explain this?
    \medskip
    \item Change \texttt{max\_features} from 500 to 50. Does the ranking of most-similar
          pairs change? What does this tell you about the sensitivity of TF-IDF to vocabulary truncation?
  \end{enumerate}
  \medskip
  \small A stretch extension comparing TF-IDF with averaged GloVe embeddings appears in the appendix.
\end{frame}

% ---------------------------------------------------------------
\section{Discussion}

\begin{frame}{Discussion: When BoW vs.\ Embeddings?}
  \begin{block}{The central trade-off}
    TF-IDF is interpretable, fast, and requires no pre-training.\\
    Word embeddings are dense, semantically rich, but are black-box vectors.
  \end{block}
  \medskip
  \textbf{Discuss in groups (5 min):}
  \begin{enumerate}
    \item You are building a system to flag 10-K filings that are unusually negative.
          Would you use TF-IDF + LM dictionary or word embeddings? Justify.
    \item You are trying to retrieve, from a database of 500,000 earnings call transcripts,
          all calls that discuss supply-chain risk. Which method would you prefer?
    \item An auditor wants a system that can \emph{explain} why a particular filing was
          flagged. Which method is easier to explain?
    \item Under what conditions might TF-IDF actually \emph{outperform} embeddings in
          a return prediction task?
  \end{enumerate}
\end{frame}

% ---------------------------------------------------------------
\section{Solutions}

\begin{frame}{Solution: Problem 1 (TF-IDF)}
  \textbf{Document frequencies and IDF} ($N=4$, $\ln$ base):
  \medskip
  \begin{center}
  \scriptsize
  \begin{tabular}{@{}lrrl@{}}
    \toprule
    Token & df & idf & $\mathrm{tfidf}$ in $d_1$ (tf = 0.2) \\
    \midrule
    declined    & 1 & $\ln 4 \approx 1.386$ & 0 \\
    growth      & 2 & $\ln 2 \approx 0.693$ & $0.2 \times 0.693 \approx 0.139$ \\
    guidance    & 1 & $1.386$ & 0 \\
    increased   & 1 & $1.386$ & $0.2 \times 1.386 \approx 0.277$ \\
    margin      & 2 & $0.693$ & $0.2 \times 0.693 \approx 0.139$ \\
    pressure    & 2 & $0.693$ & 0 \\
    revenue     & 2 & $0.693$ & $0.2 \times 0.693 \approx 0.139$ \\
    strong      & 2 & $0.693$ & $0.2 \times 0.693 \approx 0.139$ \\
    uncertainty & 2 & $0.693$ & 0 \\
    volume      & 2 & $0.693$ & 0 \\
    \bottomrule
  \end{tabular}
  \end{center}
  \medskip
  $\mathrm{sim}(d_1, d_2) \approx 0.21$ (shared: \texttt{margin}).
  Most similar pair: $d_2$ and $d_4$ (shared: \texttt{uncertainty, pressure, growth}).
\end{frame}

\begin{frame}{Solution: Problem 2 (Analogies) and Session Wrap-Up}
  \textbf{Analogy solutions:}
  \begin{itemize}
    \item \texttt{bond} $-$ \texttt{coupon} $+$ \texttt{dividend} $\approx$ \texttt{equity}\\
          {\small Bond pays coupons; equity pays dividends. Subtracting the income-type distinguishes the instrument.}
    \item \texttt{long} $-$ \texttt{profit} $+$ \texttt{loss} $\approx$ \texttt{short}\\
          {\small Long profits from price increases; short profits from price decreases, i.e.\ losses for a long.}
    \item \texttt{inflation} $-$ \texttt{high} $+$ \texttt{low} $\approx$ \texttt{deflation}\\
          {\small Directional dimension in the macro vocabulary.}
  \end{itemize}
  \medskip
  \textbf{Why these work:} co-occurrence patterns encode systematic financial relationships.
  Words that appear in the same contexts as ``bond'' but not in the contexts of ``coupon''
  tend to be equity-related terms.
  \medskip
  \begin{block}{Session Summary}
    TF-IDF: fast, interpretable, order-invariant.\\
    Word2Vec/GloVe: semantic geometry, linear analogies, domain-specific adaptation needed.\\
    Next lecture: Transformers learn \emph{context-dependent} embeddings --- no static vectors.
  \end{block}
\end{frame}

% ---------------------------------------------------------------
\appendixstart{Stretch Problem: TF-IDF vs.\ GloVe Embeddings}

\begin{frame}[noframenumbering]{Stretch: TF-IDF vs.\ Averaged GloVe Embeddings}
  \textbf{Extension (optional):} Load \texttt{glove-wiki-gigaword-100} via \texttt{gensim}.
  For each document, compute the average word embedding.
  \medskip
  \begin{enumerate}
    \item Does the embedding-based similarity ranking agree with the TF-IDF ranking?
    \item What does \emph{disagreement} between the two rankings tell you about
          what each representation captures?
    \item When would you trust the embedding ranking over the TF-IDF ranking,
          and vice versa?
  \end{enumerate}
\end{frame}

\end{document}
